\begin{figure}[p]
\centering
\begin{tikzpicture}[>=Latex,font=\small,
 card/.style={draw=blue!40!black,rounded corners=3pt,fill=blue!3,align=left,text width=5.65cm,minimum height=2.3cm,inner sep=9pt},
 arr/.style={->,thick,draw=blue!55!black}]
\node[card] (a) at (0,0) {\textbf{1. A gap in inherited reasoning}\\[3pt]
Lower degree $n(R+1)$; the matching upper certificate needs a large flattening.\\[3pt]\footnotesize v001: 8 nodes, 0 new inferences};
\node[card] (b) at (7.1,0) {\textbf{2. Boundary certificates}\\[3pt]
Symmetry makes complete tests feasible. Exact degrees 9 and 12 suggest a uniform target.\\[3pt]\footnotesize v002: 19 nodes, 7 inferences};
\node[card] (c) at (0,-3.4) {\textbf{3. A uniform pencil theorem}\\[3pt]
Adjugates and a commutator give\\$3m+3\le\delta_m\le3m+6$.\\[3pt]\footnotesize v003: 43 nodes, 22 inferences};
\node[card] (d) at (7.1,-3.4) {\textbf{4. Retain the remaining question}\\[3pt]
Multiple-block tests preserve the exact endpoint as open. The base identity remains reusable.\\[3pt]\footnotesize v004: 46 nodes, 22 inferences};
\node[card,fill=green!5] (e) at (0,-6.8) {\textbf{5. Charge contracted modes}\\[3pt]
Transfer costs the factor $n/k$. The degree-9 base and geometry give $\delta_2(T)=3n$.\\[3pt]\footnotesize v005: 66 nodes, 33 inferences};
\node[card,fill=green!5] (f) at (7.1,-6.8) {\textbf{6. Publish and preserve}\\[3pt]
Audit the original criteria; save proofs, checks, graph paths, and further targets.\\[3pt]\footnotesize v006: 68 nodes, 34 inferences};
\draw[arr] (a)--(b);
\draw[arr] (b.south)--($(b.south)!0.5!(d.north)$)-| (c.north);
\draw[arr] (c)--(d);
\draw[arr] (d.south)--($(d.south)!0.5!(f.north)$)-| (e.north);
\draw[arr] (e)--(f);
\end{tikzpicture}
\caption{Recorded research progression. Each stage changes what the next investigation can use. Arrows describe exploration history, while the theorem's implication structure is shown separately in \cref{fig:certificate-proof}. The open pencil endpoint survives the success of the two-term branch.}
\label{fig:discovery-timeline}
\end{figure}
